% Appendix C. Regime detail supporting §6 and §7. One table.

\section{The fixed book across regimes}\label{sec:appendix_c}

Table~\ref{tab:regimes} cuts the fixed book's daily output into five market
regimes. Because the book is sized once on 2019 and never changed, every
difference between rows is the market rather than the book, which makes this
the cleanest view of how a static portfolio behaves through a regime change.

% Generated by build/make_tables.py from regime_table.csv. Do not edit.

\begin{table}[tbp]
\centering
\footnotesize
\setlength{\tabcolsep}{4.5pt}
\caption{The fixed book across market regimes}
\label{tab:regimes}
\begin{tabular}{@{}lrrrrrrr@{}}
\toprule
 & & Mean & Diversif. & Realised & Exception & Equity vol & 10y vol \\
Regime & Days & VaR & (\%) & P\&L s.d. & rate (\%) & (ann., \%) & (bp/day) \\
\midrule
Calm, 2017--2019 & 687 & 54.8 & 44.8 & 28.7 & 3.5 & 12.7 & 4.0 \\
COVID shock, 2020 & 82 & 56.8 & 45.4 & 67.9 & 11.0 & 57.2 & 8.8 \\
Recovery, 2021 & 231 & 52.9 & 45.8 & 31.0 & 4.8 & 13.7 & 4.5 \\
Rate shock, 2022 & 195 & 64.9 & 42.3 & 94.9 & 32.3 & 25.4 & 8.4 \\
After the shock, 2024--2025 & 452 & 99.7 & 32.1 & 49.2 & 3.3 & 16.7 & 5.8 \\
\bottomrule
\end{tabular}
\notetext{The book is sized once on 2019 and never changed, so every difference across rows is the market rather than the book. Mean VaR and realised profit-and-loss standard deviation are in millions of US dollars at the book's normalising scale and carry no information in levels; their ratio within a row does. The last row is where the downward drift in the measured diversification effect shows up, with VaR at its sample peak and volatility back to ordinary levels.}
\end{table}


Three rows carry the argument of Section~\ref{sec:assumptions}.

The COVID window shows the lag. Equity volatility ran at 57.2\% annualised and
realised profit and loss had a standard deviation of 67.9, while VaR averaged
56.8 --- barely above its calm-period level of 54.8 --- because the shock was
six weeks old and the five-year scenario set was still overwhelmingly quiet
data. Roughly one day in nine breached.

The 2022 rate shock shows the same mechanism at its worst, and it is the
clearest cost of running the disclosed window without the disclosed weighting.
The exception rate reaches 32.3\%. Five years of scenarios in early 2022 reach
back to 2017, a stretch containing almost no rate volatility, so the engine
entered the fastest tightening cycle in four decades carrying a scenario set
that had never seen one.

The final row is where the downward drift in the measured diversification
effect appears. Mean VaR sits at its sample peak of 99.7 with volatility back
to ordinary levels and the exception rate collapsed to 3.3\%, and the
diversification effect has fallen to 32.1\% from 44.8\% in the calm years. The
book did not change. The window did: it now contains 2022 and will continue to
until 2027. That ratchet is what Section~\ref{sec:results} identifies as the
source of the trend in the model series, and it is the behaviour a dealer
repositioning daily does not exhibit.
